\documentclass[10pt,letterpaper]{article}
\usepackage{cloudsuite}
\renewcommand{\DocShort}{CH. 11}
\begin{document}
\coverpage{Chapter 11: Cloud Security and Cybersecurity Data-Oriented Mechanisms}{Cornell study notes and applied review | MECHANISMS}{Concepts \textbullet{} decisions \textbullet{} verification artifacts}
\sourceStatement
\tableofcontents
\newpage
\section{Learning contract}
\begin{takeaway}
Build data lifecycle safeguards with observable recovery outcomes.
\end{takeaway}
\textbf{Chapter scope:} Malware controls, DLP, platform trust, backup and monitoring.\\
\sourceRef{11}{311}
\subsection{Mapped source sections}
\begin{itemize}
\item 11.1 Digital Virus Scanning and Decryption System
\item 11.2 Malicious Code Analysis System
\item 11.3 Data Loss Prevention (DLP) System
\item 11.4 Trusted Platform Module (TPM)
\item 11.5 Data Backup and Recovery System
\item 11.6 Activity Log Monitor
\item 11.7 Traffic Monitor
\item 11.8 Data Loss Protection Monitor
\end{itemize}
\section{Cornell cue and notes matrix}
\small\setlength{\LTleft}{0pt}\setlength{\LTright}{0pt}
\begin{longtable}{>{\RaggedRight\arraybackslash}p{0.29\textwidth}>{\RaggedRight\arraybackslash}p{0.63\textwidth}}
\toprule\textbf{Cue / essential question}&\textbf{Explanatory notes and reasoning}\\\midrule\endfirsthead\toprule\textbf{Cue}&\textbf{Notes}\\\midrule\endhead
\textbf{How do malware controls protect data?} & Scanning and analysis can identify harmful content, but encrypted formats, false positives and execution isolation complicate complete coverage.\\[3pt]\textcolor{petrol}{\emph{Practice:} Propose a quarantine workflow with evidence handling.}\\[8pt]\midrule
\textbf{What can DLP do?} & Data loss prevention policies attempt to detect or prevent sensitive information crossing disallowed boundaries. Classification and exception governance are prerequisites.\\[3pt]\textcolor{petrol}{\emph{Practice:} Create a test dataset with permitted and blocked flows.}\\[8pt]\midrule
\textbf{Where does a TPM contribute?} & Hardware-assisted keys and measurement can anchor platform trust; the specific assurance depends on provisioning and attestation design.\\[3pt]\textcolor{petrol}{\emph{Practice:} Specify one attestation-related acceptance test.}\\[8pt]\midrule
\textbf{What makes backup a recovery control?} & A successful copy is not the same as successful restoration. Recovery objectives, isolation and restore drills determine operational confidence.\\[3pt]\textcolor{petrol}{\emph{Practice:} Record observed recovery time and the oldest recovered transaction.}\\[8pt]\midrule
\textbf{Why use activity and traffic monitors?} & Logs help reconstruct actions and traffic measurement exposes abnormal flows or performance. Their coverage must align with the incident model.\\[3pt]\textcolor{petrol}{\emph{Practice:} Build a minimum viable log event schema.}\\[8pt]\midrule
\textbf{How are data-loss monitors used?} & Monitors supply detection evidence that should lead to triage, disposition and changes in protective controls.\\[3pt]\textcolor{petrol}{\emph{Practice:} Define a closure record for a simulated incident.}\\[8pt]\midrule
\end{longtable}\normalsize
\section{Synthesis and architecture connections}
Treat each mechanism as a deployable capability, not a complete architecture. Identify the infrastructure prerequisite, management-plane owner, dependent controls and the negative/failure evidence expected from an implementation.
\begin{workbox}[title={Applied deliverable}]
Build data lifecycle safeguards with observable recovery outcomes.\\[4pt]Create a one-page architecture or operations artifact with: (1) assumptions, (2) a decision and alternatives, (3) control ownership, (4) observable acceptance criteria, and (5) a failed-case test.
\end{workbox}
\section{Active recall}
\begin{enumerate}
\item How do malware controls protect data?
\item What can DLP do?
\item Where does a TPM contribute?
\item What makes backup a recovery control?
\end{enumerate}
\subsection{Answer-check rubric}
\begin{itemize}
\item A correct response distinguishes the named concepts rather than relying on vendor names.
\item The explanation identifies an operational consequence, boundary or tradeoff.
\item At least one measurement, test or documentary artifact could demonstrate the claim.
\end{itemize}
\section{Blank Cornell capture sheet}
\textbf{Topic / source subsection:}\enspace\rule{0.65\linewidth}{0.4pt}\\[6pt]
\begin{longtable}{p{0.28\textwidth}p{0.64\textwidth}}
\toprule Cues / questions & Notes / evidence \\ \midrule
\rule{0pt}{42pt} & \\ \midrule
\rule{0pt}{42pt} & \\ \midrule
\rule{0pt}{42pt} & \\ \midrule
\rule{0pt}{42pt} & \\ \midrule
\rule{0pt}{42pt} & \\ \midrule
\end{longtable}
\textbf{Summary in three sentences:}\\[16pt]\hrulefill\\[16pt]\hrulefill\\[16pt]\hrulefill

\refs
\end{document}
